\documentclass[11pt,a4paper]{article}
\usepackage[margin=26mm,headheight=15pt]{geometry}
\usepackage[T1]{fontenc}
\usepackage[utf8]{inputenc}
\usepackage{lmodern,microtype}
\usepackage{amsmath,amssymb,amsthm,mathtools}
\usepackage{booktabs,array,longtable}
\usepackage{enumitem}
\usepackage{xurl}
\usepackage[hidelinks]{hyperref}
\usepackage{fancyhdr}
\usepackage{listings}
\hypersetup{pdftitle={Support-Controlled Reversion and the Exact One-Exponential Substitution Group},pdfauthor={OpenAI, prepared for Vladimir Reshetnikov},pdfsubject={Hahn power-log series, Lagrange inversion, Newton reversion and exponential sectors}}
\pagestyle{fancy}\fancyhf{}
\fancyhead[L]{\small Support-controlled transseries reversion}
\fancyhead[R]{\small Research article}
\fancyfoot[C]{\thepage}
\setlength{\parskip}{0.25em}
\setlength{\parindent}{1.2em}
\setlist{itemsep=0.25em,topsep=0.4em}
\allowdisplaybreaks
\emergencystretch=2em
\newtheorem{theorem}{Theorem}[section]
\newtheorem{proposition}[theorem]{Proposition}
\newtheorem{lemma}[theorem]{Lemma}
\newtheorem{corollary}[theorem]{Corollary}
\theoremstyle{definition}
\newtheorem{definition}[theorem]{Definition}
\newtheorem{example}[theorem]{Example}
\newtheorem{question}[theorem]{Research question}
\theoremstyle{remark}
\newtheorem{remark}[theorem]{Remark}
\newcommand{\R}{\mathbb R}
\newcommand{\Q}{\mathbb Q}
\newcommand{\Z}{\mathbb Z}
\newcommand{\N}{\mathbb N}
\newcommand{\C}{\mathcal C_\Gamma}
\newcommand{\Clf}{\mathcal C_{\Gamma,\mathrm{lf}}}
\newcommand{\E}{\mathcal E_{\Gamma,\mu}}
\newcommand{\supp}{\operatorname{supp}}
\newcommand{\lc}{\operatorname{lc}_t}
\newcommand{\ordq}{\operatorname{ord}_q}
\newcommand{\id}{\operatorname{id}}
\newcommand{\rev}{^{\langle-1\rangle}}
\newcommand{\coef}[1]{\left[#1\right]}
\newcommand{\OO}{\mathcal O}
\newcommand{\D}{\mathsf D}
\newcommand{\repoLabel}[1]{\texttt{\detokenize{#1}}}
\title{\textbf{Support-Controlled Reversion and the Exact\ One-Exponential Substitution Group}\\[0.65em]
\large Dense Hahn exponents, Newton denominator stability,\\
Cayley sector coefficients, and a convergent mixed example}
\author{Prepared for Vladimir Reshetnikov\\\small Research development based on the ProveIt transseries corpus}
\date{29 September 2026}
\begin{document}
\maketitle
\begin{abstract}
We address three explicitly stated questions in the ProveIt volume on
polynomial--logarithmic transseries. First, dense real exponent groups do not
require a new index set for Lagrange--B\"urmann reversion: the ordinary
positive-integer sum remains valid, with a precise valuation bound and a
support certificate. Second, every Newton iterate stays in the input
Puiseux lattice; more strongly, its normalized support remains in the positive
monoid generated by the normalized input support. Third, adjoining a single
exponential does not preserve the entire polynomial--logarithmic substitution
group. We give an exact obstruction and a replacement theorem. In
$\R[L]((t^\Gamma))[[q]]$, with $t=X^{-1}$, $L=\log X$ and
$q=e^{-\mu X}$, a map $X+P(L)+U+h$, where $v(U)>0$ and $h\in q\E$, admits the
normalized differential-compatible substitution precisely when
$P(L)=aL+b$ and $\mu a\in\Gamma$. This class is maximal among maps with a sublinear Hahn base displacement.
The admissible maps form a group, and inversion is given by an explicitly summable formula after the zero-sector
core is inverted exactly. For a single exponential perturbation we determine
the leading block of every inverse sector: its coefficient is governed by
$n^{n-1}/n!$, with exact exponent and logarithmic degree. A worked example,
$X+\log X+e^{-X}$, has rational sector coefficients, a sharp pole hierarchy,
and a uniform geometric remainder bound on a positive real ray.

All central assertions are proved here. These are resolutions of specified
repository questions and explicit extensions of its calculus, not a claim
that the classical theory of transseries reversion is new. In particular,
formal exponential-sector closure is not Borel summability or resurgence.
The accompanying finite symbolic checks are not a Lean formalization.
\end{abstract}
\newpage
\tableofcontents
\newpage

\section{Research target, provenance, and scope}
\label{sec:scope}

\subsection{The questions being answered}
The canonical ProveIt source is
\emph{Transseries: the polynomial--logarithmic calculus, series reversal at
infinity, and the inversion of rapidly growing functions}, in the
\texttt{Transseries\_And\_Inversion} package under the Fabius research
frontiers \cite{ProveIt,ProveItReadme}. Its extension chapter asks for a
Lagrange--B\"urmann formula at dense exponents, a bound on Puiseux denominators
in Newton iteration, and a one-exponential enlargement with workable
composition and reversion. We use the source's stable labels, rather than
page numbers, to specify the targets:
\begin{center}
\begin{tabular}{@{}p{0.47\textwidth}p{0.46\textwidth}@{}}
\toprule
Repository question & Result of this article\\
\midrule
\repoLabel{plt:rmk:ext-open-burmann} &
Integer-indexed formula on the full Hahn class, with an exact finite
index bound for each coefficient; Section~\ref{sec:LB}.\\[0.5em]
\repoLabel{plt:rmk:ext-open-denominators} &
$\sigma(s)=s$ suffices and is minimal under divisibility; a stronger
support-monoid statement holds; Section~\ref{sec:newton}.\\[0.5em]
\repoLabel{plt:rmk:ext-open-resurgence} &
A negative answer to unrestricted ``verbatim'' closure, an exact
admissibility criterion, and a closed formal sector calculus;
Sections~\ref{sec:sector}--\ref{sec:core}.\\
\bottomrule
\end{tabular}
\end{center}
The last row is deliberately qualified. The repository question includes
analytic language. We settle its formal algebraic component and identify an
obstruction to the proposed unrestricted substitution class. We do not
construct a universal resurgent completion or prove analytic realizability
for arbitrary formal inputs.

The source snapshot inspected through the canonical raw-source address and
the package README was retrieved on 29 September 2026. Web caching and ongoing
repository edits can make displayed snapshots differ; consequently no
unverified commit-to-render parity is asserted. The exact labels above make
the mathematical comparison reproducible. The inspection was targeted at
these statements and their prerequisites, not an exhaustive audit of every
file in the repository.

\subsection{What is classical, and what is developed here}
The background mechanisms are classical: Hahn summability, formal Lagrange
inversion, filtered fixed points, and formal substitution. The standard
transseries literature already develops composition and inversion in much
larger settings \cite{vdH,EdgarComposition,EdgarGrids}. Gessel gives a detailed
account of Lagrange inversion \cite{Gessel}. More recent work treats Taylor
expansions for generalized series and systematic correspondences between
summable derivations and automorphisms
\cite{BagayokoMantova,BKKPS,BagayokoLie}. None of these general subjects is
claimed as a new discovery here.

The contribution is a precise resolution of the listed questions in the
repository's chosen coefficient algebra, with explicit support bounds, a
necessary-and-sufficient fixed-alphabet criterion, and an all-sector leading
coefficient theorem. The proofs also distinguish three notions that must not
be conflated: existence of a formal coefficient, finiteness of an entire
truncation, and analytic convergence after evaluation. No literature-priority
claim is made for the individual formulas or for the classification in a
possibly different existing notation.

We work over $\R$ to make all real exponents and the constants $e^{b}$
unambiguous. The purely power--logarithmic results work over any
characteristic-zero field containing the exponent group. Extending constants
is not the issue addressed here.

\section{A support calculus for Hahn power--log series}
\label{sec:hahn}

\subsection{The coefficient ring and its two kinds of finiteness}
Fix an additive subgroup $\Gamma\subseteq\R$ containing $1$. Write
\[
 X=t^{-1},\qquad L=\log X
\]
as formal notation. Define $\C$ to consist of all expressions
\begin{equation}\label{eq:C}
 f=\sum_{\gamma\in S}p_\gamma(L)t^\gamma,
 \qquad p_\gamma\in\R[L],
\end{equation}
where $S\subseteq\Gamma$ is well ordered in the usual order and zero
coefficients are omitted. The notation $\R[L]((t^\Gamma))$ will mean this
\emph{Hahn ring}, not a ring restricted to a discrete lattice. Each block is a
polynomial in $L$, but its degree can grow without a uniform bound.
For $f\ne0$ set
\[
 v(f)=\min\supp f,\qquad \lc(f)=p_{v(f)}(L),\qquad v(0)=+\infty.
\]
The subring $\Clf$ consists of the \emph{left-finite} elements: for every
$B\in\R$, only finitely many support exponents are at most $B$.

A family $(f_i)_{i\in I}$ is \emph{Hahn summable} if its union of supports is
well ordered and, for every $\gamma$, only finitely many $i$ have a nonzero
$t^\gamma$ block. Its sum is then defined blockwise. This definition does not
sum infinitely many real contributions to a single coefficient, even when
those contributions converge numerically.

\begin{lemma}[Support lemma]\label{lem:support}
If $S,T\subseteq\R$ are well ordered, then $S+T$ is well ordered and every
real number has only finitely many representations $s+t$ with
$(s,t)\in S\times T$. If $S\subseteq(0,\infty)$ is nonempty and well ordered,
and $\delta=\min S$, its additive monoid
\[
 M(S)=\{0\}\cup\bigcup_{r\ge1}(\underbrace{S+\cdots+S}_{r\text{ times}})
\]
is well ordered. Every target has only finitely many representations as an
ordered word in $S$. If $S$ is left finite, then $M(S)$ is left finite.
\end{lemma}
\begin{proof}
An infinite family of different representations of one target would contain
an increasing sequence of first summands and hence a decreasing sequence of
second summands, impossible in a well-ordered set. If $S+T$ had a strictly
decreasing sequence, choose one representation of each term. A sequence in a
well-ordered set has an infinite nondecreasing subsequence; along such a
subsequence of first summands, the second summands would strictly decrease.
This proves the first assertion and, by induction, its finite-sum analogue.

A word of length $r$ in positive elements of $S$ has sum at least $r\delta$.
Below any fixed upper bound, only finitely many lengths can therefore occur.
For these finitely many lengths the preceding assertions apply. A nonempty
subset of $M(S)$ has a least element by first restricting below any one of
its elements and taking a minimum in a finite union of well-ordered sets.
The same length bound gives finiteness of the ordered words for a prescribed
sum. In the left-finite case, a word with sum at most $B$ uses only the
finitely many letters of $S\cap(0,B]$, and has bounded length. There are only
finitely many such words.
\end{proof}

\begin{lemma}[Elementary algebra and summation]\label{lem:algebra}
The Cauchy product makes $\C$ a commutative integral domain, and
$v(fg)=v(f)+v(g)$. An element $f\ne0$ is a unit if and only if $\lc(f)$ is a
nonzero real constant. If $v(f_j)\to+\infty$, the sequence-family $(f_j)$ is
Hahn summable. Every sequence whose successive errors have valuation tending
to $+\infty$ converges in the valuation topology. These statements restrict
to $\Clf$ when all the supports involved are left finite.
\end{lemma}
\begin{proof}
The support lemma makes each product coefficient finite. The leading block
of a product is the product of the leading blocks, proving the valuation
identity. An inverse forces those leading blocks to be units of $\R[L]$,
which are precisely the nonzero constants. Conversely, factor a candidate
unit as $ct^\gamma(1+u)$ with $v(u)>0$ and use
$\sum_{j\ge0}(-u)^j$, whose summability follows from the support lemma.

If the valuations of a family indexed by $j\in\N$ tend to infinity, below
any fixed bound only finitely many of its supports occur. Finite unions are
well ordered, and each coefficient has a finite fiber. This also proves the
usual Cauchy-sequence assertion by stabilized initial segments, or by a
subsequence with successively increasing error valuations. Left finiteness
is preserved because at a fixed bound there are only finitely many input
supports, each with finitely many relevant exponents.
\end{proof}

\subsection{Differentiation and near-identity substitution}
Define a strong derivation by
\begin{equation}\label{eq:D0}
 D_0(t^\gamma p(L))=t^{\gamma+1}\bigl(p'(L)-\gamma p(L)\bigr).
\end{equation}
It obeys $D_0X=1$, $D_0L=t$, and
\begin{equation}\label{eq:Dgain}
 v(D_0^j f)\ge v(f)+j.
\end{equation}
Here and below an inequality involving $v(0)$ has the usual extended-value
meaning. Repeated application gives the block identity
\begin{equation}\label{eq:blockD}
 D_0^j\bigl(t^\beta p(L)\bigr)
 =t^{\beta+j}\prod_{i=0}^{j-1}(\partial_L-\beta-i)p(L).
\end{equation}
An empty product of operators is the identity.

For $B\in\C$ with $v(B)>-1$, put $u=tB$. Thus $v(u)>0$. Define
substitution by $H=X+B$ on the generators by
\begin{equation}\label{eq:gens}
 t^\gamma\circ H=t^\gamma(1+u)^{-\gamma},\qquad
 L\circ H=L+\log(1+u).
\end{equation}
Binomial and logarithmic series in $u$ are Hahn summable. For a general
$f\in\C$ this gives the useful equivalent formula
\begin{equation}\label{eq:Taylor}
 f\circ(X+B)=\sum_{j\ge0}\frac{B^j}{j!}D_0^j f.
\end{equation}
The summand has valuation at least $v(f)+j(1+v(B))$, so the outer sum is
summable, including for supports with infinitely many elements in a bounded
interval. For each fixed $j$, the support lemma justifies its products and
its termwise differentiation.

\begin{proposition}[Substitution estimates]\label{prop:subst}
The rule \eqref{eq:Taylor} is a strongly summation-preserving algebra
homomorphism. It preserves the leading block and the valuation of nonzero
$f$. It satisfies the chain rule and associative composition. If
$v(B),v(C)\ge\alpha>-1$, then
\begin{equation}\label{eq:diffbound}
 v\bigl(f\circ(X+B)-f\circ(X+C)\bigr)
 \ge v(B-C)+v(f)+1.
\end{equation}
The maps $X+B$ with $v(B)>-1$ form a group under composition.
\end{proposition}
\begin{proof}
The $j\ge1$ terms in \eqref{eq:Taylor} have strictly larger valuation than
$f$, proving preservation of its leading block. The Leibniz rule and the
binomial identity prove multiplicativity after collecting summable terms.
The identities in \eqref{eq:gens}, and the identities for logarithms and
powers of $1+u$, prove associativity on each monomial and logarithmic block.
The same identities differentiated prove the chain rule. To pass to Hahn
sums, note that the union of input supports has a common lower bound, so the
Taylor index is bounded for each target exponent; the support lemma then
proves strong summability and permits the rearrangements.

For the difference estimate use
$B^j-C^j=(B-C)\sum_{i=0}^{j-1}B^iC^{j-1-i}$. Its contribution at index $j\ge1$
has valuation at least
\[
 v(B-C)+v(f)+1+(j-1)(\alpha+1),
\]
which proves \eqref{eq:diffbound}. Inverting $X+A$ amounts to solving
$B=-A\circ(X+B)$. This map preserves $v(B)\ge v(A)=\alpha$ and improves
difference valuation by at least $\alpha+1>0$. Iteration from $0$ therefore
converges by Lemma~\ref{lem:algebra}; the same bound proves uniqueness.
Associativity gives two-sided inversion: constructing a right inverse for
the right inverse and associating the two identities shows that the first
right inverse is also a left inverse. Closure under composition follows
directly from preservation of valuation and the formula
$(X+A)\circ(X+B)=X+B+A\circ(X+B)$.
\end{proof}

\begin{remark}[The Archimedean hypothesis has a job]
We use $\Gamma\subseteq\R$, not an arbitrary ordered abelian group. For every
$\delta>0$, the multiples $j\delta$ are cofinal in $\R$. That fact justifies
the bounds on the Taylor and iteration indices. Replacing $\Gamma$ by a
higher-rank ordered group without changing the argument would create a gap.
\end{remark}

\section{Dense exponents: the ordinary Lagrange sum is enough}
\label{sec:LB}

\subsection{A formal coefficient lemma}
We record the algebraic form of classical Lagrange inversion to isolate
exactly where the index set comes from \cite{Gessel}.

\begin{lemma}[Auxiliary-parameter inversion]\label{lem:aux}
Let $R$ be a commutative $\Q$-algebra, $\phi(z),\psi(z)\in R[[z]]$, and let
$U(s)\in sR[[s]]$ be the unique solution of $U=s\phi(U)$. For $n\ge1$,
\begin{equation}\label{eq:aux}
 [s^n]\psi(U)=\frac1n[z^{n-1}]\psi'(z)\phi(z)^n.
\end{equation}
No invertibility hypothesis on $\phi(0)$ is necessary.
\end{lemma}
\begin{proof}
Existence and uniqueness follow coefficient by coefficient. At a fixed
order both sides of \eqref{eq:aux} are universal polynomials, over $\Q$, in
finitely many coefficients of $\phi$ and $\psi$. It is therefore enough to
work in the universal polynomial ring with $\phi(0)$ inverted, where
$s=z/\phi(z)$ is an invertible change of local parameter. The formal residue
change of variables gives
\begin{align*}
 [s^n]\psi(U)
 &=\mathop{\rm Res}_{z=0}
 \psi(z)\left(\frac{z}{\phi(z)}\right)^{-n-1}
 \left(\frac{z}{\phi(z)}\right)' dz\\
 &=\mathop{\rm Res}_{z=0}
 \psi(z)\left(z^{-n-1}\phi^n-z^{-n}\phi^{n-1}\phi'\right)dz\\
 &=\frac1n\mathop{\rm Res}_{z=0}z^{-n}\psi'(z)\phi(z)^n dz.
\end{align*}
The last step uses the zero residue of the derivative of
$\psi(z)\phi(z)^n z^{-n}$. The universal polynomial ring injects into its
localization, so the polynomial identity already holds before inversion.
Specialization proves it for the original $R$.
\end{proof}

\subsection{The summable formula and its coefficient bound}
\begin{theorem}[Hahn Lagrange--B\"urmann reversion]\label{thm:LB}
Let $A\ne0$ belong to $\C$, put $\alpha=v(A)>-1$ and
$\delta=\alpha+1$, and let $F=X+A$. Then
\begin{equation}\label{eq:LB}
 F\rev=X+B,\qquad
 B=\sum_{r\ge1}\frac{(-1)^r}{r!}D_0^{r-1}(A^r).
\end{equation}
This is a Hahn-summable, positive-integer-indexed family. Its $r$th term has
valuation at least $r\delta-1$. Consequently the coefficient of $t^\lambda$
receives contributions only from
\begin{equation}\label{eq:rbound}
 1\le r\le\left\lfloor\frac{\lambda+1}{\delta}\right\rfloor.
\end{equation}
Writing $A=\sum p_\gamma t^\gamma$, that coefficient is explicitly
\begin{equation}\label{eq:blockLB}
 [t^\lambda]B=
 \sum_{1\le r\le\lfloor(\lambda+1)/\delta\rfloor}
 \frac{(-1)^r}{r!}
 \left\{\prod_{j=0}^{r-2}
       (\partial_L-\lambda+r-1-j)\right\}
 [t^{\lambda-r+1}]A^r.
\end{equation}
The formula remains valid when the support of $A$ is not left finite.
\end{theorem}
\begin{proof}
Apply Lemma~\ref{lem:aux} in $R=\C$, using the formal Taylor series
\[
 \phi(z)=-\sum_{j\ge0}\frac{D_0^jA}{j!}z^j,\qquad \psi(z)=z.
\]
The Taylor map to $R[[z]]$ is an algebra homomorphism. Thus
\[
 [s^r]U(s)=\frac{(-1)^r}{r!}D_0^{r-1}(A^r).
\]
This first proves the identity with the auxiliary parameter $s$; it makes no
claim about evaluating an arbitrary formal power series at $s=1$.

For these particular coefficients, \eqref{eq:Dgain} gives
\[
 v\bigl(D_0^{r-1}(A^r)\bigr)
 \ge r\alpha+r-1=r\delta-1\longrightarrow+\infty.
\]
Therefore $U(1)$ is a legitimate Hahn sum. The identity
$U(s)=-sA\circ(X+U(s))$ may be specialized to $s=1$: in a term with $j$
factors of $U$, the normalized valuation is bounded below by $\delta$ times
the sum of their positive auxiliary indices. Only finitely many such total
indices occur at a fixed target exponent, and every remaining multiplication
has finite fibers by Lemma~\ref{lem:support}. Thus all required
rearrangements are summable. The specialized equation is precisely the
inverse equation. Proposition~\ref{prop:subst} proves uniqueness and
two-sided inversion. Formula \eqref{eq:blockLB} follows from
\eqref{eq:blockD}, with $\beta=\lambda-r+1$.
\end{proof}

\begin{corollary}[A support certificate]\label{cor:monoid}
Let $S=\supp(A)+1\subseteq(0,\infty)$ and $M=M(S)$. Then
\begin{equation}\label{eq:Bsupp}
 \supp(B)+1\subseteq M\setminus\{0\}.
\end{equation}
If $A\in\Clf$, then $B\in\Clf$. To determine $B$ below an exclusive cutoff
$T$, it is enough to keep indices
$r<(T+1)/\delta$ in \eqref{eq:LB}.
\end{corollary}
\begin{proof}
The support of the $r$th term of \eqref{eq:LB} is contained in
$r\supp(A)+(r-1)=rS-1$. Differentiation can delete a block but cannot create
an exponent other than the displayed shift. Lemma~\ref{lem:support} now
proves the assertions.
\end{proof}

\begin{remark}[What the repository's dense-exponent question conflates]
The outer summation index in Lagrange inversion counts applications of a
perturbation; it does not enumerate the exponent group. Under the
repository's $v(A)\ge0$ hypothesis, $r\le\lambda+1$ remains a valid bound
for a coefficient at any real exponent $\lambda$. Density does not force a
transfinite or noninteger perturbation index. What may fail is the
\emph{finiteness of all the coefficients below a cutoff}: this fails for
non-left-finite supports, not for the integer-indexed operator formula.
\end{remark}

\subsection{Examples separating the issues}
\begin{example}[Irrational exponents]
Let $\rho>0$ be irrational and $A=t^\rho$, in
$\Gamma=\Z+\rho\Z$. The group is dense, but reversion starts
\begin{align*}
 F\rev={}&X-t^\rho-\rho t^{2\rho+1}
 -\frac{\rho(3\rho+1)}2t^{3\rho+2}\\
 &-\frac{\rho(4\rho+1)(2\rho+1)}3t^{4\rho+3}+\cdots.
\end{align*}
Every exponent lies in the single ray $r(\rho+1)-1$. A dense ambient group
need not produce a complicated support in an individual computation.
\end{example}

\begin{example}[An infinite bounded initial segment]
In $\Gamma=\Q$, take
\[
 A=\sum_{n\ge2}t^{1-1/n}.
\]
Its support is increasing and well ordered, with minimum $1/2$ and limit
$1$. It is not left finite. Theorem~\ref{thm:LB} still applies, but the
truncation of $B$ below $1$ is $-A$ and contains infinitely many blocks.
Indeed every $r\ge2$ term has valuation at least $2$. Finite outer index
range is not a finite-data algorithm for an arbitrary Hahn input.
\end{example}

\begin{example}[A sublinear but unbounded perturbation]
For $A=t^{-1/2}=\sqrt X$, the theorem applies with $\delta=1/2$. The inverse
is the familiar positive algebraic branch
\[
 F\rev=X+\frac12-\frac12\sqrt{1+4X}
 =X-X^{1/2}+\frac12-\frac18X^{-1/2}
   +\frac1{128}X^{-3/2}-\cdots.
\]
Here a fixed coefficient can require more Lagrange indices than the
integer-scale bound $r\le N+1$. The correct quantity is $v(tA)=\delta$.
\end{example}

\section{Newton reversion preserves the support monoid}
\label{sec:newton}

Let $F=X+A$ satisfy Theorem~\ref{thm:LB}, and put $H_*=F\rev$. Define
\begin{equation}\label{eq:newton}
 H_0=X,\qquad
 H_{k+1}=H_k-\frac{F\circ H_k-X}{F'\circ H_k},\qquad F'=D_0F.
\end{equation}
Every denominator is a unit with leading block $1$ because
$v(A')\ge\alpha+1=\delta>0$.

\begin{theorem}[No exponent-lattice enlargement]\label{thm:newton}
With $S=\supp(A)+1$ and $M=M(S)$, every iterate of \eqref{eq:newton} is
well defined in $\C$ and satisfies
\begin{equation}\label{eq:newtonsupp}
 \supp\bigl(t(H_k-X)\bigr)\subseteq M\setminus\{0\}.
\end{equation}
If $e_k=v(H_k-H_*)$ is finite, then
\begin{equation}\label{eq:newtonrate}
 e_{k+1}\ge 2e_k+\alpha+2,
 \qquad
 e_k\ge(2^{k+1}-1)\delta-1.
\end{equation}
The iterates converge to $H_*$ in the valuation topology.
\end{theorem}
\begin{proof}
First prove the support assertion, which is stronger than membership in
$\Gamma$. Write $H=X+t^{-1}V$, where $\supp V\subseteq M\setminus\{0\}$.
For one input block $t^\gamma p_\gamma(L)$,
\[
 t\bigl(t^\gamma p_\gamma(L)\circ H\bigr)
 =t^{\gamma+1}(1+V)^{-\gamma}
 p_\gamma\bigl(L+\log(1+V)\bigr).
\]
The support on the right is contained in $(\gamma+1)+M$.
Thus $tA\circ H$ and $t(F\circ H-X)=V+tA\circ H$ have support in
$M\setminus\{0\}$. Similarly $A'\circ H$ has support there. The inverse of
$1+A'\circ H$ is a geometric series supported in $M$, so the normalized
Newton correction has support in $M\setminus\{0\}$. Induction starts with
$V=0$. Since $\min(M\setminus\{0\})=\delta$, all iterates satisfy
$v(H_k-X)\ge\alpha$.

For the rate, put $E=H-H_*$ and $e=v(E)\ge\alpha$. Taylor-expand $F$ at
$H$ with displacement $-E$:
\[
 0=F\circ H-X-E(F'\circ H)
   +\sum_{j\ge2}\frac{(-E)^j}{j!}(F^{(j)}\circ H).
\]
For $j\ge2$, $F^{(j)}=D_0^jA$ has valuation at least $\alpha+j$;
composition preserves it. The last sum has valuation at least
$\alpha+2e+2$, since $e+1\ge\delta>0$. Dividing by the slope unit gives the
first inequality in \eqref{eq:newtonrate}. Also
$v(H_0-H_*)=v(B)=\alpha$, because the first Lagrange term is $-A$ and all
others have strictly higher valuation. Writing $p_k=e_k+1$ gives
$p_{k+1}\ge2p_k+\delta$ and $p_0=\delta$, proving the second bound.
An exact iterate remains exact; otherwise these bounds tend to infinity.
\end{proof}

\begin{corollary}[The Puiseux denominator question]\label{cor:puiseux}
For $s\in\N_{>0}$ and $A\in\R[L][[t^{1/s}]]$, every Newton iterate and
the inverse belong to $\R[L]((t^{1/s}))$. Thus the repository question has
answer $\sigma(s)=s$. Any integer-valued universal choice of $\sigma(s)$
must be divisible by $s$.
\end{corollary}
\begin{proof}
Here $S\subseteq s^{-1}\Z$ and $1\in s^{-1}\Z$, so
\eqref{eq:newtonsupp} and \eqref{eq:Bsupp} prove the inclusion. For minimality,
take $A=t^{1/s}$. Its inverse starts with $X-t^{1/s}$, and even the first
Newton iterate has that correction as its leading block. The exponent
$1/s$ belongs to $\sigma^{-1}\Z$ only if $s$ divides $\sigma$.
\end{proof}

\begin{remark}[Why incompleteness of the Puiseux union is irrelevant here]
A Cauchy sequence with input denominators tending to infinity may leave the
Puiseux union, as the repository correctly explains. But a Newton orbit
from one fixed input is not an arbitrary sequence in that union. Its
support is trapped in the invariant monoid above. The division step is by a
unit whose geometric inverse cannot introduce a new exponent denominator.
\end{remark}

\section{A complete one-exponential sector algebra}
\label{sec:sector}

\subsection{Definition, summation, and topology}
Fix $\mu>0$ and introduce the formal symbol $q$, with the intended
interpretation $q=e^{-\mu X}$. Set
\begin{equation}\label{eq:sectorRing}
 \E=\C[[q]],\qquad
 \D=D_0-\mu q\frac{\partial}{\partial q}.
\end{equation}
Thus
\begin{equation}\label{eq:sectorD}
 \D\left(\sum_{n\ge0}q^nf_n\right)
 =\sum_{n\ge0}q^n(D_0-n\mu)f_n.
\end{equation}
There is no uniform lower bound on the valuations $v(f_n)$ as $n$ varies.
For instance, $\sum_{n\ge0}q^nt^{-n^2}$ is an element of $\E$.

A family in $\E$ is called \emph{sectorwise Hahn summable} when its
coefficient family at each $q^n$ is Hahn summable in $\C$. Give $\E$ the
product topology of the valuation topologies of its coefficient copies of
$\C$. A neighborhood can impose finitely many conditions
$v(f_n)>B_n$. This topology lets both increasingly small power--log terms
in a fixed sector and tails with increasingly large sector index tend to
zero. It must not be identified with the $q$-adic topology with discrete
coefficients, or with a lexicographic valuation topology. We use $q$-adic
order for exact sector estimates as well; $q$-adic convergence implies
convergence in the product topology.

The product topology alone does not describe every Hahn sum in a
non-left-finite coefficient ring. Sectorwise Hahn summability remains the
formal summation rule. Keeping these two structures explicit avoids a
hidden analytic convergence assumption.

\begin{proposition}[Arithmetic and completeness]\label{prop:sectoralg}
The ring $\E$ is a complete topological differential algebra for the
coefficientwise valuation topology. Its multiplication and derivation
respect sectorwise Hahn summation. An element $f=\sum q^nf_n$ is a unit
precisely when $f_0$ is a unit of $\C$, equivalently when $f_0\ne0$ and
$\lc(f_0)\in\R^\times$. Exponentials and logarithms of $q$-adically small
elements are defined by their usual power series. More generally,
$\exp(g)$ is well defined if $v(g_0)>0$, by separating $g_0$ and the
positive-sector tail.
\end{proposition}
\begin{proof}
A coefficient of a product is a finite sum over sector pairs $i+j=n$.
The support lemma applies within each pair. Completeness is coordinatewise
completeness of $\C$. Formula \eqref{eq:sectorD} proves continuity and
compatibility with summation. A unit must have invertible constant
$q$-coefficient; conversely factor $f=f_0(1+h)$ with $h\in q\E$ and use the
$q$-adically summable geometric series. The exponential statements follow
from $\ordq(h^r)\ge r$ for $h\in q\E$, and from the positive valuation of
$g_0$. For general $g$ with $v(g_0)>0$, define
$\exp(g)=\exp(g_0)\exp(g-g_0)$.
\end{proof}

\begin{remark}[Why one fixed scalar weight does not suffice]
Suppose a proposed weight assigns $q^nt^\gamma$ the number $an+b\gamma$
with $b>0$. The support of $\sum q^nt^{-n^2}$ has weights
$an-bn^2$ unbounded below. Thus no fixed scalar weight of this kind describes
all of $\E$ as a bounded-below series algebra. The product/sectorwise
construction does not impose an artificial affine growth restriction on
negative power orders. A smaller, weight-bounded algebra can be useful for
algorithms, but is a different object.
\end{remark}

\subsection{Antiderivatives do not require further exponential scales}
Although $\D q=-\mu q$ does not improve the power order, the positive-sector
differential operators are particularly easy to invert.

\begin{theorem}[Sectorwise integration]\label{thm:integration}
The derivation $\D:\E\to\E$ is surjective and has kernel $\R$.
For every $n\ge1$, the operator $D_0-n\mu$ on $\C$ is bijective, with
\begin{equation}\label{eq:sectorprimitive}
 (D_0-n\mu)^{-1}f
 =-\sum_{j\ge0}\frac{D_0^j f}{(n\mu)^{j+1}}.
\end{equation}
The series is Hahn summable and restricts to the left-finite coefficient
class. The zero-sector antiderivative is the usual power--log block
antiderivative, unique up to a real constant.
\end{theorem}
\begin{proof}
The summands in \eqref{eq:sectorprimitive} have valuations at least
$v(f)+j$, so the series is summable. Multiplication by $D_0-n\mu$ telescopes,
leaving $f$. The operator is injective because, for $g\ne0$, the term
$-n\mu g$ has valuation $v(g)$ whereas $D_0g$ has strictly higher valuation.

For the zero sector, integrate a block $t^\gamma p(L)$ by seeking
$t^{\gamma-1}Q(L)$. The equation is
\[
 Q'-(\gamma-1)Q=p.
\]
If $\gamma\ne1$, its unique polynomial solution is
\[
 Q=-\sum_{j=0}^{\deg p}\frac{p^{(j)}}{(\gamma-1)^{j+1}}.
\]
At $\gamma=1$, take the polynomial primitive of $p$, with zero constant.
The supports are shifted by $-1$, hence remain well ordered. The only
kernel element occurs at exponent $0$ and is a real constant. Applying this
argument independently in every sector proves the theorem.
\end{proof}

\section{The exact fixed-alphabet substitution criterion}
\label{sec:classification}

Consider a prospective change of variable
\begin{equation}\label{eq:Gform}
 G=X+P(L)+U+h,\qquad P\in\R[L],\quad v(U)>0,\quad h\in q\E.
\end{equation}
We allow $U=0$ with $v(0)=+\infty$. Substitution of $G$ into any base element
of $\C$ is already meaningful: first substitute
$g_0=X+P(L)+U$, then Taylor-expand in $h$. At a fixed output sector only
finitely many Taylor powers of $h$ occur.

What remains is the image of $q$. The interpretation $q=e^{-\mu X}$
requires it to obey
\begin{equation}\label{eq:qchain}
 \D(q\circ G)=-\mu(\D G)(q\circ G).
\end{equation}
We require a nonzero image. This rules out the algebra homomorphism that
simply kills every positive sector, which is not a change of variable.

\begin{theorem}[Exact one-exponential admissibility]\label{thm:admissible}
For $G$ as in \eqref{eq:Gform}, there is a strongly
sectorwise-summation-preserving substitution homomorphism on $\E$ extending
the base substitution, with nonzero image of $q$ and the chain rule
\[
 \D(f\circ G)=(\D G)\,((\D f)\circ G),
\]
if and only if
\begin{equation}\label{eq:admissibility}
 P(L)=aL+b\quad\hbox{and}\quad m:=\mu a\in\Gamma.
\end{equation}
When \eqref{eq:admissibility} holds, the normalized substitution is
\begin{equation}\label{eq:qsub}
 q\circ G=e^{-\mu b}q\,t^m\exp\bigl(-\mu(U+h)\bigr).
\end{equation}
The normalization means that the leading block of its $q$-coefficient is
$e^{-\mu b}t^m$. It makes the substitution unique. Without normalization,
the image of $q$ can be multiplied by an arbitrary nonzero real constant.
\end{theorem}
\begin{proof}[Proof of necessity]
Let $W=q\circ G\ne0$, and let $k=\ordq W\ge0$. Write
$W=q^k u+\OO(q^{k+1})$, $u\in\C\setminus\{0\}$. The leading sector of
\eqref{eq:qchain} gives
\begin{equation}\label{eq:leadingqODE}
 D_0u=\mu(k-1)u-\mu\bigl(tP'(L)+D_0U\bigr)u.
\end{equation}
If $k\ne1$, the right side has valuation $v(u)$ whereas the left side has
valuation at least $v(u)+1$, a contradiction. Thus $k=1$.
Set $V=u\exp(\mu U)$. Its defining exponential is available because
$v(U)>0$, and \eqref{eq:leadingqODE} becomes
\[
 D_0V=-\mu tP'(L)V.
\]
Write its leading block as $t^\gamma p(L)$ with $\gamma\in\Gamma$ and
$p\ne0$. Comparing the block at exponent $\gamma+1$ yields
\begin{equation}\label{eq:polyobstruction}
 p'=(\gamma-\mu P')p.
\end{equation}
If $\gamma-\mu P'$ were a nonzero polynomial, its product with $p$ would
have degree at least $\deg p$, while $p'$ has degree strictly smaller or
vanishes. Hence $\gamma-\mu P'=0$ and $p'=0$. Therefore $P=aL+b$ and
$\gamma=\mu a\in\Gamma$. The same differential equation and the kernel
calculation for $D_0$ imply $V=ct^\gamma$ for a nonzero real $c$.
\end{proof}
\begin{proof}[Proof of sufficiency and uniqueness]
Assume \eqref{eq:admissibility}. The right side of \eqref{eq:qsub} is a
well-defined element of $q\E$ with an invertible leading sector
coefficient. Extend base substitution by
\begin{equation}\label{eq:fullsub}
 \left(\sum_{n\ge0}q^nf_n\right)\circ G
   =\sum_{n\ge0}(q\circ G)^n(f_n\circ G).
\end{equation}
In sector $N$ only $n\le N$ can occur; for each $n$, Taylor expansion of
$f_n$ around $g_0$ uses only finitely many powers of $h$. Within each such
term Proposition~\ref{prop:subst} preserves Hahn summability. This proves
well-definedness, multiplicativity, and strong compatibility with sectorwise
Hahn sums. Direct differentiation of \eqref{eq:qsub} proves
\eqref{eq:qchain}; the base chain rule then extends to every element.

For uniqueness, compare any other nonzero solution $W$ of
\eqref{eq:qchain} with the canonical one $W_0$. Necessity proved that both
have sector order $1$ and that the coefficient of $q$ is a unit of $\C$.
Thus $W/W_0$ belongs to $\E$ (divide after canceling $q$), and its derivative
is zero. Theorem~\ref{thm:integration} implies that this quotient is a real
constant. Normalization makes that constant $1$.
\end{proof}

\begin{corollary}[A genuine obstruction, not a convergence defect]
For the ordinary integer-power algebra $\R[L]((t))[[q]]$, the map
$X+a\log X+b+U+h$ is admissible precisely when $\mu a\in\Z$.
The map $X+L^2$ is not admissible for any additive exponent subgroup
$\Gamma\subseteq\R$. Neither changing the topology on the same algebra nor
allowing all real power exponents removes the latter obstruction.
\end{corollary}
\begin{proof}
The first statement specializes Theorem~\ref{thm:admissible}. For $P=L^2$,
\eqref{eq:polyobstruction} is impossible for every $\gamma$. This is a
polynomial differential equation obstruction inside the coefficient ring,
not a failure to choose sufficiently fine truncations. Formally the missing
factor is $e^{-\mu L^2}$, which is not a power--logarithmic block series of
the prescribed kind.
\end{proof}

\begin{corollary}[Minimal power-group enlargement]\label{cor:enlargement}
To permit a specified finite set of affine logarithmic shifts with
coefficients $a_1,\ldots,a_r$, the smallest additive power group containing
the original $\Gamma$ is
\[
 \Gamma'=\Gamma+\Z\mu a_1+\cdots+\Z\mu a_r.
\]
The same assertion holds for any specified set of shifts by taking its
additive span. No enlargement of the power group alone permits a polynomial
logarithmic shift of degree at least two.
\end{corollary}
\begin{proof}
Necessity follows by applying the criterion separately to each shift.
Sufficiency follows by constructing the algebra over $\Gamma'$ and applying
the theorem. The final assertion was just proved.
\end{proof}

\begin{corollary}[Maximality among sublinear base displacements]
\label{cor:maximal}
More generally, let $G=X+A+h$ with $A\in\C$, $v(A)>-1$, and $h\in q\E$.
It admits a substitution of the kind in Theorem~\ref{thm:admissible} if and
only if
\[
 A=aL+b+U,\qquad \mu a\in\Gamma,\quad v(U)>0.
\]
In particular, a power displacement such as $X+\sqrt X$ cannot be admitted
in this fixed one-exponential algebra, even when $1/2\in\Gamma$.
\end{corollary}
\begin{proof}
Repeat the leading-sector argument with $D_0A$ in place of
$tP'+D_0U$. Since $v(D_0A)>0$, the same comparison first forces
$\ordq(q\circ G)=1$. Its leading coefficient $u\ne0$ then satisfies
$D_0u=-\mu(D_0A)u$. If $\alpha=v(A)<0$, the leading block of $D_0A$
cannot vanish: its polynomial coefficient is $p'-\alpha p$ for a nonzero
polynomial $p$, and $\alpha\ne0$. Thus
\[
 v((D_0A)u)=v(u)+\alpha+1<v(u)+1\le v(D_0u),
\]
a contradiction. Hence $v(A)\ge0$, and $A=P(L)+U$ with $v(U)>0$.
Theorem~\ref{thm:admissible} finishes both directions. The case $A=0$
is included by taking $a=b=U=0$.
\end{proof}

\section{A substitution group and a core-normalized inverse}
\label{sec:core}

\subsection{The group structure}
Let $\mathcal G_{\Gamma,\mu}$ be the set of maps
\begin{equation}\label{eq:admissiblegroup}
 G=X+aL+b+U+h,\quad
 \mu a\in\Gamma,\quad v(U)>0,\quad h\in q\E,
\end{equation}
with the normalized substitution of Theorem~\ref{thm:admissible}.
Let $\mathcal G_0$ be the same class with $h=0$, and
$\mathcal G_q=\{X+h:h\in q\E\}$.

\begin{theorem}[Closed substitution group]\label{thm:group}
The class $\mathcal G_{\Gamma,\mu}$ is a group under composition. Its
substitutions are continuous automorphisms of $\E$ for the product topology.
Reduction to sector zero gives a split exact sequence
\begin{equation}\label{eq:split}
 1\longrightarrow\mathcal G_q\longrightarrow\mathcal G_{\Gamma,\mu}
 \longrightarrow\mathcal G_0\longrightarrow1.
\end{equation}
In particular, this is a semidirect product of the positive-sector group
and the admissible base group. The pair $(a,b)$ is additive under
composition and changes sign under inversion.
\end{theorem}
\begin{proof}
First, the base near-identity group exists by
Proposition~\ref{prop:subst}. Composing two base maps adds their
zero-power polynomials $aL+b$; all changes to $L$ are of positive power
valuation. Thus $\mathcal G_0$ is closed under composition and its inverses
have the form $X-aL-b+V$, $v(V)>0$. The lattice condition is preserved.

For $h\in q\E$, full substitution by $X+h$ is simply
\begin{equation}\label{eq:qTaylor}
 f\circ(X+h)=\sum_{j\ge0}\frac{h^j}{j!}\D^jf.
\end{equation}
The derivation preserves sector order, so the $j$th term has sector order
at least $\ordq f+j$. Taylor identities prove associativity and the chain
rule with exact finite algebra in each output sector. The equation for the
inverse correction $b=-h\circ(X+b)$ is a $q$-adic contraction: its difference
at two inputs has sector order at least one greater than the difference of
the inputs. Its unique solution gives a two-sided inverse as in
Proposition~\ref{prop:subst}. Hence $\mathcal G_q$ is a group.

Every map $G=g_0+h$ in \eqref{eq:admissiblegroup} factors uniquely as
\[
 G=(X+k)\circ g_0,\qquad k=h\circ g_0\rev\in q\E.
\]
The right-hand side has the canonical substitution of
Theorem~\ref{thm:admissible}: on base elements this is ordinary Taylor
substitution, and on $q$ the exponential factors multiply to
\eqref{eq:qsub}. Equivalently, the chain rule and the normalized leading
$q$-block determine that image uniquely. This verifies associativity across
the two classes. Sector-zero reduction is a homomorphism with kernel
$\mathcal G_q$ and with splitting the inclusion of $\mathcal G_0$.
Conjugating a positive-sector map by a base map again reduces to $X$ in
sector zero, so the kernel is normal. This proves \eqref{eq:split} and
inversion of every admissible map.

For continuity, fix one output sector. Only finitely many input sectors
enter \eqref{eq:fullsub}; in each of these, only finitely many powers of the
positive-sector displacement enter. Base substitution preserves valuation,
differentiation has a lower valuation bound, and multiplication by a fixed
coefficient shifts valuation by a fixed amount. A sufficiently high input
cutoff therefore implies any prescribed output cutoff. The same argument
applies to the inverse substitution.
\end{proof}

\subsection{Why the unnormalized Lagrange formula can fail}
\begin{example}[A constant shift creates infinitely many identical monomials]
Take $F=X+b+q$ with $b\ne0$. In the naive expression
\[
 \sum_{r\ge1}\frac{(-1)^r}{r!}\D^{r-1}(b+q)^r,
\]
the coefficient of $q$ in the $r$th summand is
\[
 -\frac{(\mu b)^{r-1}}{(r-1)!}.
\]
Thus the same monomial $q$ occurs at every $r$. This is not a sectorwise
Hahn-summable family. Its real numerical sum happens to converge, but that
is a different summation operation. The underlying map is nonetheless
admissible and invertible. The obstruction is to using the original
power-order proof unchanged, not to the existence of the inverse.
\end{example}

The cure is to invert the entire zero sector before expanding in the
exponential tail. This is also the principle of the repository's
exact-core inversion architecture; what must be supplied in the present
algebra is summability and closure of every substitution.

\begin{theorem}[Core-normalized Lagrange--B\"urmann formula]\label{thm:coreLB}
Let
\[
 F=F_0+E_+,\qquad F_0\in\mathcal G_0,\quad E_+\in q\E,
 \qquad H_0=F_0\rev,\quad K=E_+\circ H_0.
\]
Then $K\in q\E$ and
\begin{equation}\label{eq:coreLB}
 F\rev=H_0+
 \sum_{r\ge1}\frac{(-1)^r}{r!}
 \D^{r-1}\bigl(K^r\,D_0H_0\bigr).
\end{equation}
The family is $q$-adically summable. At sector $q^N$, only $r\le N$ can
contribute. If $\ordq E_+=d>0$, the sharper bound is $r\le\lfloor N/d\rfloor$.
\end{theorem}
\begin{proof}
Since $H_0$ is admissible, $K$ exists and has the same positive sector order
as $E_+$. The identity $F\circ H_0=X+K$ implies
$F\rev=H_0\circ(X+K)\rev$. Apply Lemma~\ref{lem:aux} with the Taylor map
for $\D$, with $\phi(z)=-K(X+z)$ and $\psi(z)=H_0(X+z)$. The general
coefficient formula gives
\[
 \frac{(-1)^r}{r!}\D^{r-1}(K^r\D H_0).
\]
Here $\D H_0=D_0H_0$. Because $\D$ preserves sector order, the displayed
coefficient has sector order at least $rd$. Specialization of the auxiliary
parameter at $1$ is therefore valid by exact finiteness in each sector.
This proves \eqref{eq:coreLB} without summing any infinite family into a
single real coefficient.
\end{proof}

\begin{example}[Exact agreement with a Lambert inverse]
For $F=X+b+q$, $H_0=X-b$ and $K=e^{\mu b}q$. Formula
\eqref{eq:coreLB} becomes
\begin{equation}\label{eq:shiftW}
 F\rev=X-b-
 \sum_{n\ge1}\frac{n^{n-1}}{n!}\mu^{n-1}e^{\mu bn}q^n
 =X-b+\frac1\mu\mathcal W(-\mu e^{\mu b}q),
\end{equation}
where $\mathcal W$ denotes the formal series inverse of $z\mapsto ze^z$ at
$0$. Substituting $y=X-b+u$ into $F(y)=X$ gives
$\mu u e^{\mu u}=-\mu e^{\mu b}q$, independently checking the formula.
Here each sector is obtained from exactly one summable Lagrange term,
unlike the naive expansion in $b+q$.
\end{example}

\section{The exact leading block of every inverse sector}
\label{sec:frontier}

The core-normalized formula yields more than existence. It gives a sharp
support boundary for every positive sector of a single-exponential
perturbation.

\begin{theorem}[Cayley leading-block law]\label{thm:cayley}
Let
\[
 F=F_0+qf,\qquad
 F_0=X+aL+b+U\in\mathcal G_0,\qquad
 f=t^\beta P(L)+\text{higher power exponents},
\]
where $f\ne0$, $\beta=v(f)$ and $P=\lc(f)\ne0$.
Write $m=\mu a$ and $\kappa=\beta-m$. Then
\[
 F\rev=H_0+\sum_{n\ge1}q^n b_n,\qquad H_0=F_0\rev,
\]
has the following exact properties for every $n\ge1$:
\begin{align}
 v(b_n)&=n\kappa,\label{eq:cayleyv}\\
 [t^{n\kappa}]b_n
 &=-\frac{n^{n-1}}{n!}\mu^{n-1}e^{\mu bn}P(L)^n.
 \label{eq:cayleyblock}
\end{align}
In particular the logarithmic degree of this leading block is
$n\deg P$, and no positive sector vanishes.
\end{theorem}
\begin{proof}
The inverse base has the form
$H_0=X-aL-b+V$, $v(V)>0$. Put
\[
 qf\circ H_0=qk,
 \qquad
 k=e^{\mu b}t^{-m}e^{-\mu V}(f\circ H_0).
\]
Base substitution preserves leading blocks. Consequently
\[
 v(k)=\kappa,\qquad \lc(k)=e^{\mu b}P(L),\qquad
 D_0H_0=1+\text{positive power order}.
\]
Only $r=n$ contributes to $q^n$ in \eqref{eq:coreLB}. Using
\eqref{eq:sectorD} gives the exact coefficient formula
\begin{equation}\label{eq:bnoperator}
 b_n=\frac{(-1)^n}{n!}(D_0-n\mu)^{n-1}
       \bigl(k^nD_0H_0\bigr).
\end{equation}
For any nonzero $g\in\C$, the leading block of $(D_0-n\mu)g$ is
$-n\mu$ times that of $g$, since $D_0g$ has power order at least one higher.
Apply this observation $n-1$ times. The sign is
$(-1)^n(-1)^{n-1}=-1$, giving \eqref{eq:cayleyblock} and proving its
nonvanishing. The degree statement follows because $\R[L]$ is an integral
domain.
\end{proof}

\subsection{Interpretation and a useful change of small parameter}
The exponent pairs of the inverse satisfy $\gamma\ge n\kappa$ in sector
$n$, and the line $\gamma=n\kappa$ is attained in every sector. This is an
exact boundary, not a heuristic balance. A logarithmic shift in the base
changes the slope from $\beta$ to $\beta-\mu a$. When this number is
negative, successively deeper exponential sectors have successively larger
power prefactors; this is allowed in $\E$ and is precisely why a uniform
lower power order was not imposed.

Introduce $Q=qt^\kappa$. Then all the coefficients of
$F\rev-H_0=\sum Q^n(t^{-n\kappa}b_n)$ have nonnegative power valuation.
Their valuation-zero blocks have generating series
\begin{equation}\label{eq:gradedW}
 \sum_{n\ge1}[t^{n\kappa}]b_n\,Q^n
 =\frac1\mu\mathcal W\bigl(-\mu e^{\mu b}P(L)Q\bigr).
\end{equation}
This is an identity of formal leading-block series; it is not in general
the full inverse. The tree coefficients $n^{n-1}$ are classical
\cite{Gessel}. Their precise occurrence here follows from the leading
sector action $D_0-n\mu\sim-n\mu$, while differentiation of power--log
coefficients contributes only above the support boundary.

\begin{remark}[What the theorem does and does not generalize]
A tail with several primitive exponential sectors can have additional
contributions and cancellations at a fixed sector. Formula
\eqref{eq:coreLB} still applies, but the exact nonvanishing statement in
Theorem~\ref{thm:cayley} uses the hypothesis $E_+=qf$. Likewise, an arbitrary
formal coefficient $f$ need not be realizable by an analytic function.
Neither assumption may be silently removed.
\end{remark}

\section{A mixed logarithmic--exponential inverse with certified convergence}
\label{sec:mixed}

We now specialize to $\Gamma=\Z$, $\mu=1$, and
\begin{equation}\label{eq:mixedF}
 F(X)=X+L+q.
\end{equation}
This example exhibits all the structural features established above: the
logarithmic core changes the leading power in each exponential sector,
core inversion makes the coefficient formula finite, and the resulting
formal series has an analytic realization with a quantitative error bound.
The analytic proof is specific to this example and does not assert
convergence of arbitrary elements of $\E$.

\subsection{An exact rational coefficient formula}
Let $w=H_0(X)$ be the formal inverse of $X+\log X$. Thus
\begin{equation}\label{eq:omegaRelation}
 w+\log w=X,\qquad
 D_0w=\frac{w}{w+1},\qquad q\circ w=qw.
\end{equation}
Here $\log w=L+\log(tw)$ is legitimate because $tw=1+$ a positive-valuation
series. The last identity follows from the normalized exponential
substitution: $X-w=\log w$. Analytically, the positive solution is the
Wright omega value $w=W(e^X)$; the repository already develops this core
\cite{ProveIt,ProveItReadme}.

Define the rational differential operator
\[
 T=\frac{w}{w+1}\frac{d}{dw}.
\]
Theorem~\ref{thm:coreLB}, with $K=qw$ and $D_0H_0=w/(w+1)$, gives the
following exact formula:
\begin{equation}\label{eq:bnmixed}
 F\rev(X)=w+\sum_{n\ge1}q^n b_n(w),\qquad
 b_n(w)=\frac{(-1)^n}{n!}(T-n)^{n-1}
       \frac{w^{n+1}}{w+1}.
\end{equation}
Rational functions in this display are interpreted formally by their
expansions at $w=\infty$ and then by substitution of $w=H_0(X)$. Their only
possible finite poles will be at $w=-1$, so this interpretation is
unambiguous. The first three coefficients are
\begin{align}
 b_1(w)&=-\frac{w^2}{w+1},\label{eq:b1}\\
 b_2(w)&=-\frac{w^3(2w^2+2w-1)}{2(w+1)^3},\label{eq:b2}\\
 b_3(w)&=-\frac{w^4(9w^4+18w^3-11w+1)}{6(w+1)^5}.
 \label{eq:b3}
\end{align}
The coefficient functions have no hidden infinite sum: obtaining $b_n$
requires exactly $n-1$ applications of a rational differential operator.

It is also useful to regard the exponential parameter as an independent
complex variable $z$. If $\delta=\sum_{n\ge1}b_n(w)z^n$, its defining
identity is
\begin{equation}\label{eq:deltaImplicit}
 \delta+\log\left(1+\frac{\delta}{w}\right)
       +wz e^{-\delta}=0.
\end{equation}
To see this, substitute $w+\delta$ into \eqref{eq:mixedF}, use
\eqref{eq:omegaRelation}, and then replace $q$ by $z$. The coefficient of
$\delta$ at $(\delta,z)=(0,0)$ is $1+1/w$, so the identity also determines
all $b_n$ recursively whenever $w\ne0,-1$; for the analytic result below
we use only real $w\ge1$.

\subsection{Sharp poles and exact growth at infinity}
\begin{theorem}[Two exact coefficient boundaries]\label{thm:mixedpoles}
For every $n\ge1$, $b_n$ is rational, has no finite pole except $w=-1$,
and has a pole of exact order $2n-1$ there. With $(-1)!!=1$,
\begin{equation}\label{eq:polelimit}
 \lim_{w\to-1}(w+1)^{2n-1}b_n(w)
       =-\frac{(2n-3)!!}{n!}.
\end{equation}
At infinity,
\begin{equation}\label{eq:bninfty}
 b_n(w)=-\frac{n^{n-1}}{n!}w^n+\OO(w^{n-1}).
\end{equation}
\end{theorem}
\begin{proof}
The operator $T-n$ introduces no pole away from $w=-1$. Put $u=w+1$.
Then
\[
 T=\frac{-1+u}{u}\frac{d}{du}.
\]
Applied to a leading pole $c u^{-k}$, its most singular term is
$kc u^{-k-2}$; the other part of $T$, and multiplication by $-n$, produce
strictly smaller pole orders. The initial function
$w^{n+1}/(w+1)$ has leading pole $(-1)^{n+1}u^{-1}$.
After $n-1$ applications, the pole order is $1+2(n-1)$ and its
coefficient is
\[
 (-1)^{n+1}(1\cdot3\cdot5\cdots(2n-3)).
\]
Multiplication by $(-1)^n/n!$ proves \eqref{eq:polelimit}, which is nonzero.

For the second assertion, the initial rational function is
$w^n+\OO(w^{n-1})$. The operator $T$ lowers its leading power by one,
whereas multiplication by $-n$ preserves it. Consequently the leading
term after $n-1$ operations is $(-n)^{n-1}w^n$, and the prefactor in
\eqref{eq:bnmixed} gives \eqref{eq:bninfty}.
\end{proof}

Equation~\eqref{eq:bninfty} is the specialization of
Theorem~\ref{thm:cayley} with $a=1$, $\beta=0$, $P=1$, hence
$\kappa=-1$. The inverse's $n$th exponential sector begins at $t^{-n}$.
The exponential is still the asymptotically small quantity; its polynomial
prefactor grows with $n$.

The deepest pole pieces alone have the formal sum
\begin{equation}\label{eq:deepestResum}
 -\sum_{n\ge1}\frac{(2n-3)!!}{n!}
          \frac{z^n}{(w+1)^{2n-1}}
 =(w+1)\left(\sqrt{1-\frac{2z}{(w+1)^2}}-1\right).
\end{equation}
This identity resums only the most singular Laurent term of each $b_n$,
not the complete inverse. Analogous deepest-pole and square-root mechanisms
already occur in the repository's special-function inversion chapter,
notably \repoLabel{p6:thm:deepest-pole} \cite{ProveIt}. The point here is an
exact realization for the mixed logarithmic--exponential map and its
compatibility with the support calculation above.

\subsection{A uniform analytic sector bound}
\begin{theorem}[Convergence and explicit remainder]\label{thm:analytic}
For every real $w\ge1$, equation \eqref{eq:deltaImplicit} has a unique
solution in $|\delta|\le1/4$ for $|z|\le1/(16w)$. This solution is
holomorphic in $z$ in the interior, has Taylor coefficients $b_n(w)$ from
\eqref{eq:bnmixed}, and satisfies
\begin{equation}\label{eq:coefbound}
 |b_n(w)|\le\frac14(16w)^n\qquad(n\ge1).
\end{equation}
For $\theta=16w|z|<1$ and every integer $N\ge0$,
\begin{equation}\label{eq:uniformbound}
 \left|\delta(w,z)-\sum_{n=1}^{N}b_n(w)z^n\right|
 \le \frac{\theta^{N+1}}{4(1-\theta)}.
\end{equation}
In particular, set $x=w+\log w$ and $z=e^{-x}$. If $w\ge\log32$, then
$\theta=16e^{-w}\le1/2$, the positive real inverse of
$y\mapsto y+\log y+e^{-y}$ is
\[
 y=w+\sum_{n\ge1}b_n(w)e^{-nx},
\]
and its $N$-sector truncation has error at most
\begin{equation}\label{eq:realbound}
 \frac{(16e^{-w})^{N+1}}{4(1-16e^{-w})}
 \le\frac12(16e^{-w})^{N+1}.
\end{equation}
\end{theorem}
\begin{proof}
Use the principal logarithm on $|\delta/w|\le1/4$, put $s=1+1/w$, and
rewrite \eqref{eq:deltaImplicit} as the fixed-point equation
\begin{equation}\label{eq:contraction}
 \delta=\mathcal T_z(\delta):=
 -\frac{\log(1+\delta/w)-\delta/w+wz e^{-\delta}}{s}.
\end{equation}
For $|u|<1$,
\[
 |\log(1+u)-u|\le\frac{|u|^2}{2(1-|u|)}.
\]
Thus for $|\delta|\le1/4$, $w\ge1$, and $|z|\le1/(16w)$,
\[
 |\log(1+\delta/w)-\delta/w|\le\frac1{24},\qquad
 |wz e^{-\delta}|\le\frac{e^{1/4}}{16}\le\frac1{12}.
\]
The last inequality uses $e^{1/4}\le(1-1/4)^{-1}=4/3$. Since $s\ge1$,
we obtain $|\mathcal T_z(\delta)|\le1/8$, strictly inside the proposed
disc.

Its derivative with respect to $\delta$ satisfies
\begin{align*}
 |\partial_\delta\mathcal T_z(\delta)|
 &\le\frac1s\left(
       \left|\frac1{w+\delta}-\frac1w\right|
       +w|z|e^{1/4}\right)\\
 &\le\frac1s\left(\frac{1/4}{w(w-1/4)}+\frac1{12}\right)
 \le\frac5{12}<1.
\end{align*}
The closed disc is convex, so this derivative estimate is a contraction
bound. The contraction theorem supplies existence and uniqueness there.
Picard iteration from zero consists of holomorphic functions of $z$ and
converges uniformly, by the same contraction constant, on every compact
subdisc. Hence the fixed point is holomorphic. At $z=0$ it is zero, and
comparison of coefficients in \eqref{eq:deltaImplicit} identifies its
Taylor coefficients with the formal ones.

The bound $|\delta|\le1/4$ and Cauchy's estimate on circles of radius
$r<1/(16w)$ imply $|b_n(w)|\le(1/4)r^{-n}$. Letting $r$ increase to
$1/(16w)$ proves \eqref{eq:coefbound}. Summing the resulting geometric
majorant for $n>N$ proves \eqref{eq:uniformbound}.

For the real specialization, $we^{-x}=e^{-w}$, so the estimates give
\eqref{eq:realbound}. The fixed point is real by uniqueness and conjugation,
and $y=w+\delta>0$. Moreover,
\[
 \frac{d}{dy}(y+\log y+e^{-y})=1+\frac1y-e^{-y}>0\qquad(y>0).
\]
This map tends to $-\infty$ at $0^+$ and to $+\infty$ at infinity;
therefore its positive inverse is unique and is the solution just
constructed.
\end{proof}

The bound is intentionally simple and conservative. It is uniform in the
number of retained sectors, not an asymptotic statement with an unspecified
constant depending on $N$. It does \emph{not} solve the repository's general
uniformity question for truncated Wright omega or Lambert expansions:
$w$ is kept as an exactly inverted core, rather than replaced by its own
power--logarithmic truncation. Controlling that second truncation uniformly
requires a separate estimate.

\subsection{Reproducible numerical checks}
The package's script evaluates the first six coefficients at 100 decimal
digits and compares truncations with a directly computed root. Selected
results are shown below. These floating-point checks illustrate the proven
bound; they are not interval-arithmetic certificates and are not used in its
proof.
\begin{center}
\begin{tabular}{@{}rrrr@{}}
\toprule
$w$ & $N$ & Observed absolute error & Bound \eqref{eq:realbound}\\
\midrule
4 & 2 & $4.5901\cdot10^{-6}$ & $8.8998\cdot10^{-3}$\\
4 & 4 & $3.2414\cdot10^{-9}$ & $7.6430\cdot10^{-4}$\\
4 & 6 & $2.9452\cdot10^{-12}$ & $6.5637\cdot10^{-5}$\\
8 & 2 & $3.9227\cdot10^{-11}$ & $3.8866\cdot10^{-8}$\\
8 & 4 & $1.1932\cdot10^{-17}$ & $1.1197\cdot10^{-12}$\\
8 & 6 & $4.6841\cdot10^{-24}$ & $3.2258\cdot10^{-17}$\\
\bottomrule
\end{tabular}
\end{center}
Values in the last column have been rounded upward at the displayed
precision. The full output also includes $w=5$ and $w=12$, all orders from
one through six, and the associated $x=w+\log w$.

\section{Computation, proof certificates, and formalization}
\label{sec:computational}

\subsection{What the support bounds actually make finite}
For an input $A$ with $\delta=1+v(A)>0$, the Lagrange formula has a finite
\emph{outer} index range below any prescribed power cutoff $T$. It is enough
to retain the positive integers $r$ with
\[
 r\delta-1<T.
\]
This is not yet a finite algorithm for every Hahn input. A well-ordered
support may contain infinitely many exponents below $T$, as the example
in Section~\ref{sec:LB} shows. A finite truncation algorithm additionally
requires left-finiteness or a symbolic representation that can handle such
an infinite initial segment. This distinction is essential when stating
complexity claims.

For left-finite input, the following support certificate is sufficient:
record the positive set $S=1+\supp A$, its least element $\delta$, and the
monoid $M(S)$ as a generated object. Products, substitutions, unit inverses,
and Newton steps are then authorized by their support inclusions. One need
not form all of $M(S)$ in advance; to work below $T$, only bounded-length
words in the finitely many generators that can contribute need be visited.
The finite-factorization lemma guarantees termination. For dense rational
or irrational exponent groups, exact comparison of input exponents is an
additional representation requirement, not a consequence of the theorem.

Newton's estimate gives a second certificate. To guarantee agreement with
the inverse at every exponent strictly below $T$, it is enough to choose
$k$ so that
\[
 (2^{k+1}-1)\delta-1\ge T.
\]
The denominator invariant is independent of this choice: more iterations
do not force a finer Puiseux lattice. This separates precision growth from
exponent-representation growth.

\subsection{Sector-first inversion}
For a general admissible map, the efficient organization is:
\begin{enumerate}[label=\arabic*.]
 \item Invert the zero-sector map exactly in the formal power--logarithmic
       sense, obtaining $H_0$ with the desired certified power truncation.
 \item Transport the positive-sector perturbation to
       $K=E_+\circ H_0$ and record $d=\ordq K$.
 \item For sector $N$, apply \eqref{eq:coreLB} only for
       $1\le r\le\lfloor N/d\rfloor$.
 \item Within those finitely many terms, propagate power lower bounds
       through multiplication and differentiation before truncating.
\end{enumerate}
A uniform power cutoff applied blindly to every intermediate factor is
unsafe: multiplication by a negative-valuation coefficient can move an
omitted term into the requested range. The bounds from
Theorem~\ref{thm:cayley} give a particularly simple remedy for a single
primitive sector: normalize by $Q=qt^\kappa$, so all resulting coefficient
valuations are nonnegative. For arbitrary positive-sector tails, use
sector-dependent lower bounds instead of assuming this single straight
boundary.

Core normalization is not merely an optimization. It changes a generally
nonsummable Lagrange family, such as the constant-shift example of
Section~\ref{sec:core}, into a family with a strict increase of exponential
order. A computer algebra system should reject the unnormalized operation
unless an additional summation structure has explicitly been supplied.

\subsection{The included verification artifact}
The accompanying \texttt{verify\_results.py} has two different roles.
Its exact symbolic part checks:
\begin{itemize}
 \item mixed rational-exponent Lagrange reversion for
 $A=L+t^{1/2}(L^2+1)+t^{2/3}$ modulo terms of valuation at least $4$;
 \item agreement with three Newton steps at that precision, with all
 exponents remaining in $\tfrac16\Z$;
 \item the rational formulas for $b_1,\ldots,b_6$, their exact pole
 coefficients and leading powers, and the implicit equation
 \eqref{eq:deltaImplicit} through $z^6$.
\end{itemize}
Its numerical part produces the table data for the analytic example.
The exact part uses rational arithmetic and symbolic polynomial/rational
simplification in SymPy. The numerical part uses mpmath. A text receipt
and a CSV file are supplied so that the reported checks are inspectable
without running code. These checks are useful for catching signs,
normalizations, and low-order mistakes; they do not replace the quantified
proofs in this article.

\subsection{A realistic Lean dependency plan}
The repository README already distinguishes formalized scale, Neumann,
block, and coefficient lemmas from an assembled composition/reversion
calculus \cite{ProveItReadme}. The results in this article should be
formalized at that same statement-level granularity. The main dependency
sequence is:
\begin{center}
\begin{tabular}{@{}p{0.19\textwidth}p{0.74\textwidth}@{}}
\toprule
Layer & Required mathematical interface\\
\midrule
Support & Well-ordered support addition; finite fibers; positive-monoid
closure; valuation-complete coefficient ring.\\[0.35em]
Differentiation & A strongly summable derivation with valuation gain;
finite polynomial block identities.\\[0.35em]
Substitution & A Taylor homomorphism with quantitative gain, associativity,
chain rule, and the difference estimate.\\[0.35em]
Reversion & Auxiliary-parameter Lagrange inversion; justified specialization;
fixed-point uniqueness; Newton error and support invariants.\\[0.35em]
Exponential & Ordinary power series over the coefficient ring; sectorwise
derivation; kernel and integration; leading-block obstruction.\\[0.35em]
Core inversion & Group factorization; finite sector coefficient extraction;
Cayley boundary and worked rational coefficients.\\
\bottomrule
\end{tabular}
\end{center}
No new Lean module has been compiled for this article. In particular, a
symbolic identity verified by the included Python program is not being
promoted to a kernel-checked theorem. The benefit of the present formulation
is that both the positive results and the obstruction have small, explicit
interfaces suitable for independent formalization.

\section{Further research questions}
\label{sec:questions}

The questions below are proposed directions, not claims of established open
status in the entire literature. Some sharpen a repository frontier; others
arise from the exact boundary proved here. Their hypotheses are chosen so
that a future result can be compared with the present theorems without
silently changing the ambient algebra.

\begin{question}[The minimal enlargement for quadratic logarithmic shifts]
What is the smallest explicitly presented coefficient algebra extending
$\R[L]((t^\Gamma))$ that is closed under a specified family of maps
$X+P(L)+U$ with $\deg P\ge2$, and under their inverses? Adjoining
$e^{-\mu L^2}$ resolves one immediate missing monomial, but repeated
composition also changes its exponent through
$\log(X+P(L)+U)=L+$ smaller terms. Can an effective support description
track all descendants while keeping coefficient extraction finite?
\end{question}
Theorem~\ref{thm:admissible} proves that merely enlarging $\Gamma$ cannot
solve this problem. It does not prove that no larger transseries algebra
can do so; standard transseries fields provide much richer ambient
structures \cite{vdH,EdgarComposition}.

\begin{question}[Several exponential actions and resonances]
For $q_j=e^{-\mu_jX}$, can the admissibility and leading-sector theorems be
organized uniformly over a finitely generated action monoid, including
rational relations among the $\mu_j$? In a freely indexed model the
conditions $\mu_j a\in\Gamma$ are suggested by the one-generator theorem,
but dependent actions identify different formal multi-indices. Determine
support and noncancellation hypotheses that make the resulting coefficient
formula representation-independent.
\end{question}

\begin{question}[Analytic subclasses stable under exact-core inversion]
Give coefficient-growth and analytic continuation assumptions on a subclass
of $\E$ under which \eqref{eq:coreLB} preserves Borel summability, or a
specified resurgent structure. Which singularities can be created when the
zero-sector core is inverted, and how do the corresponding Stokes data
transform? Formal existence in the product topology alone supplies none
of these conclusions.
\end{question}

\begin{question}[Two-stage uniform error bounds]
In the mixed example, replace the exact $w$ by a truncated power--logarithmic
expansion while also truncating the exponential sectors. Find a bound uniform
in both truncation orders, with an explicit admissible relation between the
orders and $X$. More generally, formulate conditions on an invertible core
that turn the formal support estimate into an analytic residual-to-error
certificate with constants uniform over a parameter family.
\end{question}
This would connect the present result to the repository's
\repoLabel{plt:rmk:ext-open-uniformity}, without confusing a bound for the
exponential tail with a bound for the core's own asymptotic expansion.

\begin{question}[Beyond Archimedean exponent groups]
When the exponent group is an arbitrary ordered abelian group, what replaces
the real-valued hypothesis used in the coefficientwise Lagrange bound and
the sequential fixed-point proof? A positive gain $\delta$ need not have
multiples cofinal in the whole group. Identify conditions under which
strong Hahn summability still suffices, and distinguish genuine failures of
existence from failures of sequential convergence in a chosen topology.
\end{question}
The current article deliberately restricts to $\Gamma\subseteq\R$;
no non-Archimedean extension is implicit in its proofs.

\begin{question}[Exact boundaries for several primitive sectors]
If $E_+=\sum_{j=1}^{d}q^j f_j$, determine the smallest possible power
valuation at each output sector and classify when its leading coefficient
cancels. A finite minimization over sector decompositions predicts a lower
bound, but the exact boundary requires algebraic control of all minimizing
terms. Seek hypotheses under which an analogue of the nonvanishing Cayley
law can be proved.
\end{question}

\begin{question}[Realization beyond Puiseux supports]
Which left-finite rational-exponent power--logarithmic series outside every
fixed Puiseux lattice admit compatible realizations in a Hardy field, with
formal differentiation, composition, and inversion agreeing with those of
the representing germs? The support theorems here separate the formal
problem from this analytic realization problem; they do not select a
canonical germ for an arbitrary formal series.
\end{question}
This is a sharpened continuation of
\repoLabel{plt:rmk:ext-open-realization}, not a consequence of denominator
stability for a fixed input lattice.

\begin{question}[Sharp complex transition near the core critical point]
For the mixed example, can the deepest-pole expression
\eqref{eq:deepestResum} be promoted to a uniform scaling approximation near
$w=-1$, with both error bounds and the actual nearby branch points located?
The rational pole hierarchy is exact, but its formal leading-pole sum is not
yet a uniform two-variable approximation to the complete inverse.
\end{question}

\begin{question}[Proof-producing support and precision inference]
Can a symbolic engine infer the support monoid, admissible exponent lattice,
and required intermediate precisions automatically, then emit small
independently checked certificates? A useful first target is a Lean-checked
implementation of the fixed-Puiseux and finitely generated action cases,
with explicit rejection of nonsummable unnormalized substitutions.
\end{question}

\section{Conclusion}
The dense-exponent reversion problem does not require replacing the
positive-integer Lagrange index: the normalized valuation gain supplies
exactly the finiteness that density was thought to obstruct. Newton
iteration preserves more than a denominator bound---it preserves the
positive support monoid generated by the input. These conclusions settle
the stated dense-Hahn and Puiseux-iteration questions in their specified
real-exponent settings.

The one-exponential problem has a different answer. The polynomial--logarithmic
coefficient ring cannot support arbitrary logarithmic shifts once
$q=e^{-\mu X}$ must transform differentially correctly. Among sublinear Hahn base displacements, the exact surviving
class consists of affine logarithmic shifts with $\mu a\in\Gamma$, positive
power-order corrections, and arbitrary positive exponential sectors.
Within that class, block operations, integration, composition, and inversion
are rigorous and explicit. Exact-core normalization is the mechanism that
restores summable Lagrange inversion.

For a single primitive exponential perturbation, every inverse sector has
an exact leading power and a nonzero Cayley coefficient. The mixed inverse
of $X+\log X+e^{-X}$ shows how this formal statement can meet analysis:
rational coefficient functions, sharp pole orders, and a uniform geometric
remainder bound hold simultaneously. The remaining issues---richer exponent
alphabets, multi-action cancellations, analytic realization and resurgence,
and complete formal verification---are now separated from the claims
actually proved.

\begingroup
\small
\begin{thebibliography}{99}

\bibitem{ProveIt}
Vladimir Reshetnikov, ProveIt repository,
\emph{Transseries: the polynomial--logarithmic calculus, series reversal at
infinity, and the inversion of rapidly growing functions}.
Canonical source, targeted inspection on 29 September 2026:
\url{https://raw.githubusercontent.com/VladimirReshetnikov/ProveIt/main/Analysis/FabiusFunction/docs/semi-formalized-research-frontiers/drafts/series-and-transseries/Transseries_And_Inversion/transseries_and_inversion.tex}.
The stable labels cited in the text identify the relevant statements;
no current-source/current-PDF parity is assumed.

\bibitem{ProveItReadme}
ProveIt, \emph{Transseries\_And\_Inversion package README}.
Repository status and statement-level formalization inventory, inspected
29 September 2026.
\url{https://github.com/VladimirReshetnikov/ProveIt/blob/main/Analysis/FabiusFunction/docs/semi-formalized-research-frontiers/drafts/series-and-transseries/Transseries_And_Inversion/README.md}.

\bibitem{vdH}
Joris van der Hoeven, \emph{Transseries and Real Differential Algebra}.
Lecture Notes in Mathematics, vol.~1888, Springer, 2006.
\url{https://doi.org/10.1007/3-540-35590-1}.

\bibitem{Gessel}
Ira M. Gessel, Lagrange inversion.
\emph{Journal of Combinatorial Theory, Series A} \textbf{144} (2016),
212--249.
\url{https://doi.org/10.1016/j.jcta.2016.06.018}.
Preprint: \url{https://arxiv.org/abs/1609.05988}.

\bibitem{EdgarComposition}
G. A. Edgar, \emph{Transseries: Composition, Recursion, and Convergence}.
Preprint, 2009. \url{https://arxiv.org/abs/0909.1259}.

\bibitem{EdgarGrids}
G. A. Edgar, \emph{Transseries: Ratios, Grids, and Witnesses}.
Preprint, 2009. \url{https://arxiv.org/abs/0909.2430}.

\bibitem{BagayokoMantova}
Vincent Bagayoko and Vincenzo Mantova,
\emph{Taylor expansions over generalised power series}.
Preprint, 2025. \url{https://arxiv.org/abs/2509.08473}.

\bibitem{BKKPS}
Vincent Bagayoko, Lothar Sebastian Krapp, Salma Kuhlmann,
Daniel Panazzolo, and Michele Serra,
\emph{Automorphisms and derivations on algebras endowed with formal infinite
sums}. Preprint, 2024; revised version consulted.
\url{https://arxiv.org/abs/2403.05827}.

\bibitem{BagayokoLie}
Vincent Bagayoko, \emph{A formal Lie correspondence}.
Preprint, 2026. \url{https://arxiv.org/abs/2604.04224}.

\end{thebibliography}
\endgroup
\end{document}
